\section*{NeurIPS Paper Checklist}

\begin{enumerate}

\item {\bf Claims}
    \item[] Question: Do the main claims made in the abstract and introduction accurately reflect the paper's contributions and scope?
    \item[] Answer: \answerYes{}
    \item[] Justification: The abstract and introduction clearly state four contributions (taxonomy, benchmark, CoT strategies, experiments) that are all substantiated in the paper. Limitations regarding U.S.\ equities only and 7B-scale models are discussed in Section~7.

\item {\bf Limitations}
    \item[] Question: Does the paper discuss the limitations of the work performed by the authors?
    \item[] Answer: \answerYes{}
    \item[] Justification: Section~7 (Conclusion) discusses limitations including coverage of only U.S.\ equities (S\&P~500, daily closing prices), 7B-parameter models, and the performance ceiling on prediction tasks due to market efficiency.

\item {\bf Theory assumptions and proofs}
    \item[] Question: For each theoretical result, does the paper provide the full set of assumptions and a complete (and correct) proof?
    \item[] Answer: \answerNA{}
    \item[] Justification: The paper is primarily empirical and does not include theoretical results or proofs.

    \item {\bf Experimental result reproducibility}
    \item[] Question: Does the paper fully disclose all the information needed to reproduce the main experimental results of the paper to the extent that it affects the main claims and/or conclusions of the paper (regardless of whether the code and data are provided or not)?
    \item[] Answer: \answerYes{}
    \item[] Justification: Section~4 describes the training setup (LoRA hyperparameters, backbone models, epochs). Section~3 details the data construction (250 S\&P~500 stocks, splits, sample sizes). Appendix provides QA generation parameters, class distributions, and full task examples.

\item {\bf Open access to data and code}
    \item[] Question: Does the paper provide open access to the data and code, with sufficient instructions to faithfully reproduce the main experimental results, as described in supplemental material?
    \item[] Answer: \answerYes{}
    \item[] Justification: Code and data will be released upon acceptance. The paper provides sufficient detail (data construction parameters, training hyperparameters, task specifications) for independent reproduction.

\item {\bf Experimental setting/details}
    \item[] Question: Does the paper specify all the training and test details (e.g., data splits, hyperparameters, how they were chosen, type of optimizer) necessary to understand the results?
    \item[] Answer: \answerYes{}
    \item[] Justification: Training details are in Section~4 (LoRA $r{=}32$, $\alpha{=}64$, LR $5{\times}10^{-5}$, 4 epochs). Data splits are in Section~3 (Table~5). QA generation parameters are in Appendix~D.

\item {\bf Experiment statistical significance}
    \item[] Question: Does the paper report error bars suitably and correctly defined or other appropriate information about the statistical significance of the experiments?
    \item[] Answer: \answerNo{}
    \item[] Justification: We do not report error bars or confidence intervals. Instead, we evaluate on three independent test splits (Test~A/B/C) that systematically vary stock universe and time period, providing evidence of consistency across OOD conditions. Single-run results are standard practice for LLM fine-tuning at this scale.

\item {\bf Experiments compute resources}
    \item[] Question: For each experiment, does the paper provide sufficient information on the computer resources (type of compute workers, memory, time of execution) needed to reproduce the experiments?
    \item[] Answer: \answerYes{}
    \item[] Justification: Section~4 specifies 2$\times$ NVIDIA L40S GPU training. All models are 7B-parameter scale with LoRA fine-tuning.

\item {\bf Code of ethics}
    \item[] Question: Does the research conducted in the paper conform, in every respect, with the NeurIPS Code of Ethics \url{https://neurips.cc/public/EthicsGuidelines}?
    \item[] Answer: \answerYes{}
    \item[] Justification: The research uses publicly available stock price data and open-source models. No human subjects are involved. The work does not provide trading recommendations.

\item {\bf Broader impacts}
    \item[] Question: Does the paper discuss both potential positive societal impacts and negative societal impacts of the work performed?
    \item[] Answer: \answerYes{}
    \item[] Justification: The paper focuses on financial reasoning rather than trading signal generation. Section~7 notes limitations. The model's prediction accuracy (around 68\% on prediction tasks) is insufficient for reliable automated trading, and the paper explicitly notes market efficiency as a fundamental constraint.

\item {\bf Safeguards}
    \item[] Question: Does the paper describe safeguards that have been put in place for responsible release of data or models that have a high risk for misuse (e.g., pre-trained language models, image generators, or scraped datasets)?
    \item[] Answer: \answerNA{}
    \item[] Justification: The model is a 7B fine-tuned LLM for financial reasoning tasks with classification accuracy far below reliable trading thresholds. It poses minimal risk for misuse compared to general-purpose LLMs.

\item {\bf Licenses for existing assets}
    \item[] Question: Are the creators or original owners of assets (e.g., code, data, models), used in the paper, properly credited and are the license and terms of use explicitly mentioned and properly respected?
    \item[] Answer: \answerYes{}
    \item[] Justification: All baseline models (TimeOmni-1, Qwen2.5, Llama, Mistral, Gemma, Phi, TimeMQA) and forecasting models (Chronos, PatchTST, etc.) are properly cited. Stock price data is from publicly available S\&P~500 daily closing prices.

\item {\bf New assets}
    \item[] Question: Are new assets introduced in the paper well documented and is the documentation provided alongside the assets?
    \item[] Answer: \answerYes{}
    \item[] Justification: FinTSR-Bench is documented in Section~3 with task specifications (Table~4), data splits (Table~5), QA generation parameters (Appendix~D), and full task examples (Appendix~A). Code and data will be released upon acceptance.

\item {\bf Crowdsourcing and research with human subjects}
    \item[] Question: For crowdsourcing experiments and research with human subjects, does the paper include the full text of instructions given to participants and screenshots, if applicable, as well as details about compensation (if any)?
    \item[] Answer: \answerNA{}
    \item[] Justification: The paper does not involve crowdsourcing or research with human subjects. All data is programmatically generated from publicly available stock prices.

\item {\bf Institutional review board (IRB) approvals or equivalent for research with human subjects}
    \item[] Question: Does the paper describe potential risks incurred by study participants, whether such risks were disclosed to the subjects, and whether Institutional Review Board (IRB) approvals (or an equivalent approval/review based on the requirements of your country or institution) were obtained?
    \item[] Answer: \answerNA{}
    \item[] Justification: The paper does not involve human subjects research.

\item {\bf Declaration of LLM usage}
    \item[] Question: Does the paper describe the usage of LLMs if it is an important, original, or non-standard component of the core methods in this research? Note that if the LLM is used only for writing, editing, or formatting purposes and does \emph{not} impact the core methodology, scientific rigor, or originality of the research, declaration is not required.
    \item[] Answer: \answerYes{}
    \item[] Justification: LLMs are a core component of our method. We fine-tune TimeOmni-1-7B and Qwen2.5-7B-Instruct as backbone models, as described in Sections~4 and~5. All baseline LLMs used for comparison are cited with their respective references.

\end{enumerate}
